\section{The Tetrad}

As we have journeyed through each Number, they appear to us like different states of consciousness: in the 1, I see Unity, in 2, I behold the difference between Being and Non-Being, and with 3, I see the deeper harmony of the One-and-the-Many. Of course, neither Unity, Being, nor Harmony is lacking in any of the Numbers: it is just that each Number is a key with a different dominant tone. One contains all Numbers, Two is another Monad (as well as Privation), and Three reconciles mystery with mystery.

When we reach Tetrad, we find more of the same, but differently. Are these new states of Being, or just changes in Consciousness?

$1+2+3+4=10$, so the Tetrad is complete: it contains the Pyramid, the four elements, and the four causes of Aristotle, and the Quadrivium --- arithemetic apprehends quantity in general, music apprehends relative quantity, geometry takes in size in general (static), and astronomy apprehends size in motion in the stars. Iamblichus thinks the Quadrivium's use of Number to apprehend Truth is less liable to error (one supposes, than religion?). He doesn't seem to have a division between sacred/profane science -- there is only Truth. He relates each science of Quadrivium to a number: Arithmetic (Monad), Harmony/Music (Dyad), Geometry (Triad), and Astronomy (Tetrad). The first and simplest three dimensional figure is the sphere, composed of center, diameter, area, circumference.

Creation proper begins with the Tetrad, as fire is symbolized by the pyramid (4 bases, 4 angles), there are 4 elements, 4 seasons, etc. Iamblichus' use of 4 orderings in Nature is a long one. ``Accretion of discrete things cannot be seen without the Tetrad, although plurality appears with the Triad''. Those who have died in virtue are not ``thrice-blessed'', but ``four times blessed'': they have gone beyond alteration and change.

Squares symbolize stability, and Hermes and Herakles are associated with them. \emph{Tlao} means endure, so we may call it the Tetlad, for a neologism. Since its perimeter and square are equal, it is called Justice. 4 is associated with the ensoulment of Body, thanks to the various harmonic ratios established by the 1,2,3,4 sequence; so 4 (again) represents emergence out of possibility into actual existence.

In the Western tradition, we would say that the number 4 establishes the outflowing of the seminal Ideas of God into the material universe: Four is the entrance of the Triad's energies into the actual harmonies of physical creation. In Platonic terms, matter would now be Over-Souled or given an Anima, so that the world-Soul now exists, and is a living thing. The seeds of Creation or rationes seminales (associated with Logos Spermatikos) proceed to actually unify Matter and Form into substance, shape, form, \& principle. So actual things come into existence, and do not latently exist in conception as Ideas in God's mind only. God desires to save the world, as a husband desires to deliver his bride from the clutches of a dragon. This is the will of heaven, and perhaps Christianity's greatest gift to the other religions, if it is done properly and seen properly.

We should see an analogy here\footnote{\url{http://www.praxisresearch.net/custpage.cfm/frm/65464/sec\_id/65464}} with the Hindu tradition, which teaches 1) What is experienced 2) The experiencer 3) A medium of experience. 1 is Archetype (Father) 2 is the Logos, which reflects that which is received (Son) and 3 is the hypostatic union which reconciles Subject/Object. When this inner Trinity is ignited in perichoresis, Creation (4) is the result. Isn't this what is means to experience an Initiation? First one establishes consciousness of Tradition (1), then one learns obedience (primarily of mind or intuition, which alone can rule the psyche -- 2), Illumination follows (3) in which one experiences the Lord's will as one's own, and shares the reality ``in the Spirit''. At this point, the worshiper has been purgated, illumined, and experienced theosis (the term illumination in English slips back and forth between 2\&3, but technically, belongs to the hour before dawn, when the worshipper is still in actual darkness, but not formally so). Only at this point, after 3, can the devotee actually begin to ``change the world'', because their Self is changed: they can now begin to truly live, and affect Creation (4).

If we, with Tomberg, forget about our salvation, and focus on preparing in steps for what is inevitable (from the point of view of vertical energy), then we can avoid short-circuiting the process and falling victim to Counter Initiation. For some, this will come ``after death'' in the post-mortem state: according to Jesus' parables, these as fully inherit the Kingdom as those who enter the vineyard at dawn.

It is a question of seeking for what we truly want. If like Simon Magus, we desire power Now, we risk falling. If one waits upon the Lord, and trims the lamps, watching and praying, the perfect time shall come, and may come sooner, after all.


\flrightit{Posted on 2013-01-18 by Logres}

\begin{center}* * *\end{center}

\begin{footnotesize}
\begin{sffamily}
\texttt{David on 2013-01-18 at 22:26 said:}

This was a very good series of text. Thank you.

Might I ask if Iamblichus was the only one positing such ideas ? I have not read the original Pythagoreans philosophers much, but I might guess they had some ideas about mathematics and geometry regarding Truth ? If so, on what different or similar levels ?

\hfill

\texttt{Logres on 2013-01-19 at 06:26 said:}

No, they all held similar views in the sense that they all thought the same way: Number (literal) is an analogy (or more than an analogy) for higher principles. Niocomachus, etc. the whole bunch are named throughout the text and freely drawn upon : they are cited almost as authorities. It's a short text, if you ever have the desire to read it.


\hfill

\end{sffamily}
\end{footnotesize}

